\documentclass[aspectratio=169, 10pt]{beamer}

% -----------------------------------------------------------------------------
% Presentation Theme and Styling
% -----------------------------------------------------------------------------
\usetheme{Madrid}
\usecolortheme{default}
\usefonttheme{professionalfonts}

\usepackage{amsmath}
\usepackage{amssymb}
\usepackage{booktabs}
\usepackage{graphicx}
\usepackage{xcolor}
\usepackage{microtype}

% Color Definitions (Consistent with Research Overview Palette)
\definecolor{NavyBlue}{RGB}{31, 73, 125}     % #1F497D
\definecolor{TealAccent}{RGB}{0, 156, 189}    % #009CBD
\definecolor{DarkSlate}{RGB}{40, 44, 52}      % #282C34
\definecolor{CardBg}{RGB}{248, 250, 252}      % #F8FAFC
\definecolor{CardBorder}{RGB}{203, 213, 225}  % #CBD5E1
\definecolor{AlertRed}{RGB}{192, 80, 77}      % #C0504D
\definecolor{SuccessGreen}{RGB}{44, 160, 44}  % #2CA02C

% Beamer Color Assignments
\setbeamercolor{palette primary}{bg=NavyBlue, fg=white}
\setbeamercolor{palette secondary}{bg=NavyBlue!85, fg=white}
\setbeamercolor{palette tertiary}{bg=NavyBlue!70, fg=white}
\setbeamercolor{structure}{fg=NavyBlue}
\setbeamercolor{title}{fg=NavyBlue}
\setbeamercolor{frametitle}{fg=NavyBlue, bg=white}
\setbeamercolor{block title}{bg=NavyBlue!12, fg=NavyBlue}
\setbeamercolor{block body}{bg=CardBg, fg=DarkSlate}
\setbeamercolor{item}{fg=TealAccent}

% Remove default navigation symbols
\setbeamertemplate{navigation symbols}{}

% Clean footline with slide numbers
\setbeamertemplate{footline}{
  \leavevmode%
  \hbox{%
    \begin{beamercolorbox}[wd=.70\paperwidth,ht=2.5ex,dp=1.1ex,left]{author in head/foot}%
      \hspace*{2ex}\textcolor{gray}{\scriptsize Spectral Complexity Algorithm Mathematics}
    \end{beamercolorbox}%
    \begin{beamercolorbox}[wd=.30\paperwidth,ht=2.5ex,dp=1.1ex,right]{title in head/foot}%
      \textcolor{gray}{\scriptsize Slide \insertframenumber{} / \inserttotalframenumber}\hspace*{2ex}
    \end{beamercolorbox}%
  }%
  \vskip0pt%
}

% -----------------------------------------------------------------------------
% Title Information
% -----------------------------------------------------------------------------
\title[Spectral Complexity Mathematics]{Mathematical Formulation of the Spectral Complexity Algorithm}
\subtitle{Gram-Space Cholesky Maximum Distance, Simplex Volume, and Residual Noise Floor Tracking}
\author{Nick Chiaratti}
\institute{Rochester Institute of Technology \\ Chester F. Carlson Center for Imaging Science}
\date{}

\begin{document}

% -----------------------------------------------------------------------------
% Title Frame
% -----------------------------------------------------------------------------
\begin{frame}
  \titlepage
\end{frame}

% -----------------------------------------------------------------------------
% Slide 1: Geometric Simplex Volume & Orthogonal Heights
% -----------------------------------------------------------------------------
\begin{frame}{Geometric Simplex Volume from Iterative Orthogonal Heights}
  \begin{columns}[T]
    \begin{column}{0.48\textwidth}
      \begin{block}{Iterative Endmember Selection}
        \footnotesize
        \begin{itemize}\setlength{\itemsep}{1.5pt}
          \item \textbf{Anchor ($v_0$):} Initial reference pixel (e.g. furthest from origin).
          \item \textbf{Pivot 1 ($v_1$):} Max distance $h_0 = \|v_1 - v_0\|$ (1D base line).
          \item \textbf{Pivot 2 ($v_2$):} Max perpendicular height $h_1$ to segment $v_0\text{--}v_1$.
          \item \textbf{Pivot 3 ($v_3$):} Max perpendicular height $h_2$ to plane $v_0\text{--}v_1\text{--}v_2$.
        \end{itemize}
      \end{block}
    \end{column}
    
    \begin{column}{0.48\textwidth}
      \begin{block}{General Simplex Volume Formula}
        \footnotesize
        \begin{itemize}\setlength{\itemsep}{1.5pt}
          \item Simplex volume defined by product of orthogonal heights:
          $$\text{Volume}_k = \frac{1}{k!} \prod_{j=0}^{k-1} h_j = \frac{1}{k!} (h_0 \cdot h_1 \cdots h_{k-1})$$
          \item 2D Triangle Area: $\text{Vol}_2 = \frac{1}{2!} h_0 h_1 = \frac{1}{2} \cdot \text{Base} \cdot \text{Height}$
          \item 3D Tetrahedron Vol: $\text{Vol}_3 = \frac{1}{3!} h_0 h_1 h_2$
          \item Exact volume without spectral dimensionality reduction.
        \end{itemize}
      \end{block}
    \end{column}
  \end{columns}

  \vspace{0.05cm}
  \begin{center}
    \includegraphics[height=0.24\textheight]{output_hyspex_residuals/diagram_endmember_heights.png}
  \end{center}
\end{frame}

% -----------------------------------------------------------------------------
% Slide 2: Gram-Space Formulation & Initial Distances
% -----------------------------------------------------------------------------
\begin{frame}{Gram-Space Representation \& Distance Initialization}
  \begin{columns}[T]
    \begin{column}{0.48\textwidth}
      \begin{block}{1. Localized Spatial Gram Matrix}
        \footnotesize
        \begin{itemize}\setlength{\itemsep}{2pt}
          \item Local pixel window $Y = [y_1, \dots, y_N] \in \mathbb{R}^{B \times N}$ ($N=9$).
          \item Form $N \times N$ Gram matrix once per local window:
          $$K = Y^T Y \in \mathbb{R}^{N \times N}, \quad K_{ij} = y_i^T y_j$$
          \item Decouples operations from the spectral dimension $B$.
        \end{itemize}
      \end{block}

      \begin{block}{2. Distance to Initial Vertex ($v_0$)}
        \footnotesize
        \begin{itemize}\setlength{\itemsep}{2pt}
          \item Squared Euclidean distance evaluated directly from $K$:
          $$d_i^{(1)} = \|y_i - v_0\|^2 = K_{ii} - 2 K_{i, i_0} + K_{i_0, i_0}$$
          \item Requires zero access to raw $B$-dimensional spectra.
        \end{itemize}
      \end{block}
    \end{column}

    \begin{column}{0.48\textwidth}
      \begin{block}{3. Greedy Pivot Selection}
        \footnotesize
        \begin{itemize}\setlength{\itemsep}{2pt}
          \item At step $k$, select unselected pixel with max distance:
          $$i_k = \arg\max_{i \notin \mathcal{S}} d_i^{(k)}$$
          \item Orthogonal height is the square root of max distance:
          $$h_k = \sqrt{d_{i_k}^{(k)}}$$
          \item Update selected set: $\mathcal{S}_k = \mathcal{S}_{k-1} \cup \{i_k\}$.
        \end{itemize}
      \end{block}

      \begin{block}{Computational Efficiency}
        \footnotesize
        \begin{itemize}\setlength{\itemsep}{2pt}
          \item Local spatial window has $N = 9 \ll B$ (where $B=200\text{--}400$).
          \item Initializing in Gram space replaces expensive spectral vector passes with $O(N)$ lookups.
        \end{itemize}
      \end{block}
    \end{column}
  \end{columns}
\end{frame}

% -----------------------------------------------------------------------------
% Slide 3: Cholesky Factorization & Exact Schur Downdating
% -----------------------------------------------------------------------------
\begin{frame}{Cholesky Factorization \& Exact Schur Downdating}
  \begin{columns}[T]
    \begin{column}{0.48\textwidth}
      \begin{block}{4. Cholesky Factor Column Update}
        \footnotesize
        \begin{itemize}\setlength{\itemsep}{1pt}
          \item Compute column $k$ of lower triangular factor $L \in \mathbb{R}^{N \times k}$:
          $$L_{i, k} = \frac{1}{h_k} \left( K_{i, i_k} - \sum_{j=1}^{k-1} L_{i, j} L_{i_k, j} \right)$$
          \item Pivot value: $L_{i_k, k} = h_k$.
          \item Previously selected pivots: $L_{i_j, k} = 0$ for $j < k$.
        \end{itemize}
      \end{block}

      \begin{block}{5. Exact Schur Complement Downdate}
        \footnotesize
        \begin{itemize}\setlength{\itemsep}{1pt}
          \item Update squared orthogonal distances for unselected pixels:
          $$d_i^{(k+1)} = d_i^{(k)} - L_{i, k}^2$$
          \item \textbf{Subspace Equivalence:} $d_i^{(k+1)} \equiv \|P_{\mathcal{S}_k}^\perp y_i\|^2$, matching orthogonal subspace projection.
        \end{itemize}
      \end{block}
    \end{column}

    \begin{column}{0.48\textwidth}
      \begin{block}{Elimination of Matrix Inversion}
        \footnotesize
        \begin{itemize}\setlength{\itemsep}{1pt}
          \item Classic OSP explicitly constructs:
          $$P_k^\perp = I - U_k (U_k^T U_k)^{-1} U_k^T$$
          \item Incurring $O(B^2 N + B^3)$ operations at every iteration.
          \item Cholesky Schur downdating requires only $O(N \cdot k)$ operations, with zero matrix inversions.
        \end{itemize}
      \end{block}

      \begin{block}{Numerical Implementation Properties}
        \footnotesize
        \begin{itemize}\setlength{\itemsep}{1pt}
          \item Computed in full \texttt{float64} precision to eliminate loss of significance on small residuals.
          \item Yields $1.5\times\text{--}2.0\times$ acceleration over classic OSP.
          \item Preserves raw data reality without synthetic smoothing.
        \end{itemize}
      \end{block}
    \end{column}
  \end{columns}
\end{frame}

% -----------------------------------------------------------------------------
% Slide 4: Residual Tracking & Pooled Noise Variance
% -----------------------------------------------------------------------------
\begin{frame}{Residual Tracking \& Degree-of-Freedom Pooled Noise}
  \begin{columns}[T]
    \begin{column}{0.48\textwidth}
      \begin{block}{1. Total Residual Sum of Squares (RSS)}
        \footnotesize
        \begin{itemize}\setlength{\itemsep}{2pt}
          \item Summing updated distances $d_i$ across unselected pixels gives the unmodeled Frobenius energy:
          $$\text{RSS}_k = \sum_{i \notin \mathcal{S}_k} d_i^{(k+1)} \equiv \|P_{\mathcal{S}_k}^\perp Y\|_F^2$$
          \item Evaluated directly from the Schur complement downdate.
        \end{itemize}
      \end{block}

      \begin{block}{2. Degrees of Freedom Accounting}
        \footnotesize
        \begin{itemize}\setlength{\itemsep}{2pt}
          \item Total data measurements in local patch: $B \times N$.
          \item Consumes $k$ spectral directions ($B - k$ remaining).
          \item Consumes $k$ spatial pivot constraints ($N - k$ remaining).
          \item Net unmodeled degrees of freedom: $(B - k)(N - k)$.
        \end{itemize}
      \end{block}
    \end{column}

    \begin{column}{0.48\textwidth}
      \begin{block}{3. Pooled Noise Standard Deviation}
        \footnotesize
        \begin{itemize}\setlength{\itemsep}{2pt}
          \item Unbiased pooled estimate of residual noise:
          $$\hat{\sigma}_k = \sqrt{\frac{\text{RSS}_k}{(B - k)(N - k)}}$$
          \item Properly normalizes for subspace rank exhaustion.
        \end{itemize}
      \end{block}

      \begin{block}{Why Heuristic Elbow Tests Fail}
        \footnotesize
        \begin{itemize}\setlength{\itemsep}{2pt}
          \item Heuristic ratio tests ($h_k / h_1 \le 0.01$) couple to scene dynamic range $h_1 = \|v_1 - v_0\|$.
          \item Pure noise retains energy $(B-k)N\sigma_n^2$ in subspace.
          \item Greedy selection forces $h_k$ to plateau at extreme-value noise ceiling $\approx \sigma_n \sqrt{2 \ln N}$.
        \end{itemize}
      \end{block}
    \end{column}
  \end{columns}
\end{frame}

% -----------------------------------------------------------------------------
% Slide 5: Residual Variance Drop Ratio & Noise Floor Identification
% -----------------------------------------------------------------------------
\begin{frame}{Residual Variance Drop Ratio \& Noise Floor Identification}
  \begin{columns}[T]
    \begin{column}{0.48\textwidth}
      \begin{block}{Step-to-Step Variance Drop Ratio}
        \footnotesize
        \begin{itemize}\setlength{\itemsep}{2pt}
          \item Ratio of successive pooled standard deviations:
          $$\rho_k = \frac{\hat{\sigma}_{k-1}}{\hat{\sigma}_k}$$
          \item Measures fractional reduction in unmodeled variance.
        \end{itemize}
      \end{block}

      \begin{block}{Physical Material Signal Regime}
        \footnotesize
        \begin{itemize}\setlength{\itemsep}{2pt}
          \item Resolving a physical material captures substantial variance:
          $$\rho_k \ge 3.0 \quad (\text{empirically } 4.0\times\text{--}18.9\times)$$
          \item Validated on panels, boundaries, and diseased foliage.
        \end{itemize}
      \end{block}
    \end{column}

    \begin{column}{0.48\textwidth}
      \begin{block}{Isotropic Sensor Noise Floor Regime}
        \footnotesize
        \begin{itemize}\setlength{\itemsep}{2pt}
          \item When only noise remains, projecting arbitrary directions removes only nominal degrees of freedom $\approx \frac{1}{N-k}$:
          $$\rho_k \approx 1.01\text{--}1.06 \approx 1.00\times$$
          \item Collapses to theoretical isotropic noise baseline ($1.0\times$).
        \end{itemize}
      \end{block}

      \begin{block}{Scale-Free Stopping Criterion}
        \footnotesize
        \begin{itemize}\setlength{\itemsep}{2pt}
          \item Intrinsic rank $k^*$ is identified when $\rho_k$ hits noise floor:
          $$k^* = \min \{ k : \rho_{k+1} < \tau_{\text{noise}} \}, \quad \tau_{\text{noise}} \approx 1.5\times$$
          \item Operates across sensors without prior noise calibration.
        \end{itemize}
      \end{block}
    \end{column}
  \end{columns}
\end{frame}

\end{document}
